
\section{Spinor Field Lagrangian Formulation}\label{appx_spinor_lagrangian}
\subsection{Minimal Geometric Field Action}
Starting from the matrix equation, Eq.~\eqref{eq_fundamental_matrix}, on the $+c$ propagation branch, the complex column field $\Psi\in\mathbb{C}^{8}$ has the temporal-volume adjoint
\begin{equation}
    \bar{\Psi} \coloneqq \Psi^{\dagger}\Gamma^{0},
\end{equation}
and admits the minimal nongauge Lagrangian
\begin{equation}\label{eq_flag_lagrangian_sym}
    \mathcal{L}_{0} = \frac{i}{2}\bar{\Psi}\Gamma^{\mu}\overleftrightarrow{\partial}_{\mu}\Psi - \chi\bar{\Psi}R\Psi,
\end{equation}
with $R = \rho(\mathbf{R})$, $\chi = 1/\hbar$, and
\begin{equation}
    \bar{\Psi}\Gamma^{\mu}\overleftrightarrow{\partial}_{\mu}\Psi \coloneqq \bar{\Psi}\Gamma^{\mu}\partial_{\mu}\Psi - (\partial_{\mu}\bar{\Psi})\Gamma^{\mu}\Psi.
\end{equation}
Retaining the PO-source operator preserves the description of the fundamental field dynamics through the pair $(\boldsymbol{\Psi},\mathbf{R})$. For notational simplicity, $\mathbf{R}\equiv\mathbf{R}_{\mathrm{total}}$ denotes the total PO-source, whose scalar component naturally supplies the mass term as a solution to the propagation-orthogonality problem $\mathbf{U}\cdot\mathbf{R}=0$ as demonstrated below.

Treating $\Psi$ and $\bar{\Psi}$ as independent fields, stationary action yields
\begin{equation}\label{eq_flag_euler_lagrange}
    \left(i\Gamma^{\mu}\partial_{\mu} - \chi R\right)\Psi = 0,
\end{equation}
together with the adjoint equation
\begin{equation}\label{eq_flag_adjoint_euler_lagrange}
    i(\partial_{\mu}\bar{\Psi})\Gamma^{\mu} + \chi\bar{\Psi}R = 0.
\end{equation}
Up to a total divergence, Eq.~\eqref{eq_flag_lagrangian_sym} is equivalent to the unsymmetrized first-order form
\begin{equation}
    \mathcal{L}_{0} = \bar{\Psi}(i\Gamma^{\mu}\partial_{\mu} - \chi R)\Psi.
\end{equation}

\subsection{Reality of the Propagation-Orthogonal Geometric-Source}
Reality of the source bilinear gives the Hermiticity condition with respect to the temporal-volume adjoint,
\begin{equation}\label{eq_R_temporal_volume_reality}
    R^{\ddagger} = R, \quad \mathcal{O}^{\ddagger} \coloneqq \Gamma^{0}\mathcal{O}^{\dagger}\Gamma^{0}.
\end{equation}
Since $\rho(\mathbf{I}_{t})=i\Gamma^{0}$ and $R=\rho(\mathbf{R})$, this expression is the matrix representation of the abstract condition
\begin{equation}
    \mathbf{I}_{t}\mathbf{R}^{\dagger}\mathbf{I}_{t} = -\mathbf{R} \Longleftrightarrow \mathbf{e}_{0}\mathbf{R}^{\dagger}\mathbf{e}_{0} = -\mathbf{i}\mathbf{R}\mathbf{i},
\end{equation}
established in the main text.

With the stacked spinor $\Psi=(\Psi_{+},\Psi_{-})^{\mathsf{T}}$, Eq.~\eqref{eq_R_temporal_volume_reality} gives
\begin{equation}
    R_{++} = \gamma^{0}R_{++}^{\dagger}\gamma^{0}, \quad R_{--} = \gamma^{0}R_{--}^{\dagger}\gamma^{0},
\end{equation}
and consequently
\begin{equation}\label{eq_R_off_diagonal_reality}
    R_{+-} = -\gamma^{0}R_{-+}^{\dagger}\gamma^{0}, \quad R_{-+} = -\gamma^{0}R_{+-}^{\dagger}\gamma^{0}.
\end{equation}
Accordingly, the sector-preserving blocks are Hermitian with respect to the four-dimensional Dirac adjoint, whereas the Flag-time mixing blocks form anti-adjoint pairs.

\subsection{Conserved Currents and Positive Norm}
\label{ch_noether_currents}
Global phase invariance yields the total Noether current
\begin{equation}\label{eq_total_current}
    J^{\mu} = \bar{\Psi}\Gamma^{\mu}\Psi, \quad \partial_{\mu}J^{\mu} = 0,
\end{equation}
whose temporal component is
\begin{equation}
    J^{0} = \bar{\Psi}\Gamma^{0}\Psi = \Psi^{\dagger}\Psi \geq 0.
\end{equation}
Consequently, the minimal Lagrangian supplies a positive-definite norm density and assigns the same kinetic sign to both Flag-time sectors, excluding a relative wrong-sign kinetic mode.

More generally, consider any constant interaction generator $T_{a}$ satisfying
\begin{equation}\label{eq_current_generator_conditions}
    T_{a}^{\ddagger} = T_{a}, \quad [T_{a}, \Gamma^{\mu}] = 0, \quad \partial_{\mu}T_{a} = 0.
\end{equation}
For such a generator, the corresponding bilinear current is
\begin{equation}\label{eq_internal_noether_current}
    J_{a}^{\mu} \coloneqq \bar{\Psi}\Gamma^{\mu}T_{a}\Psi.
\end{equation}
Substitution of the field and adjoint equations then yields the exact divergence balance
\begin{equation}\label{eq_internal_current_exchange}
    \partial_{\mu}J_{a}^{\mu} = -i\chi\bar{\Psi}[T_{a}, R]\Psi.
\end{equation}
Taking $T=I_{8}$ recovers the total current, since the commutator vanishes identically.

Among these constant generators, those commuting with $R$ throughout the source background form its stabilizer algebra,
\begin{equation}\label{eq_R_stabilizer_algebra}
    \mathfrak{h}_{R} \coloneqq \left\{T \in \mathfrak{g} \mid [T, R] = 0\right\}.
\end{equation}
For generators outside $\mathfrak{h}_{R}$, Eq.~\eqref{eq_internal_current_exchange} describes the source-dependent exchange of the corresponding charge. In this way, the PO-source identifies the conserved field symmetries and supplies the exchange terms for the remaining currents.

\subsection{Symmetric Hyperbolicity and Causal Propagation}
Because the algebraic PO-source operator modifies neither the differential order nor the principal symbol
\begin{equation}
    \mathcal{P}(\xi) = \Gamma^{\mu}\xi_{\mu},
\end{equation}
the Clifford relations yield
\begin{equation}
    \det\mathcal{P}(\xi) = \left(\eta_{\mathrm{4D}}^{\mu\nu}\xi_{\mu}\xi_{\nu}\right)^{4},
\end{equation}
so the characteristic set remains the ordinary Lorentzian null cone.

Multiplying Eq.~\eqref{eq_flag_euler_lagrange} by $-i\Gamma^{0}$ produces a first-order system with principal matrices
\begin{equation}
    A^{0} = \Gamma^{0}\Gamma^{0} = I_{8} > 0, \quad A^{j} = \Gamma^{0}\Gamma^{j} = \left(A^{j}\right)^{\dagger}.
\end{equation}
Consequently, the field equation is symmetric hyperbolic, and for sufficiently regular PO-source coefficients its Cauchy problem is locally well posed on spacelike hypersurfaces, yielding unique solutions that depend continuously on the initial data. Since $R$ is algebraic in the spinor equation, it leaves the characteristic cone unchanged and introduces no higher derivatives of the spinor field.

\subsection{Mass, Scalar, and Pseudoscalar Interactions from the Plane-of-Now Algebra}
Within the two-dimensional Plane-of-Now, the Clifford algebra contains scalar, vector, and bivector grades. Let
\begin{equation}
    \mathbf{w}(\tau) = \cos\vartheta(\tau)\mathbf{e}_{5} + \sin\vartheta(\tau)\mathbf{e}_{6},\quad \mathbf{w}_{\perp}\coloneqq\mathbf{i}\mathbf{w}
\end{equation}
define two orthogonal unit vectors within this plane, with spacetime-dependent amplitudes evaluated along $X^{\mu}(\tau)$. The three grade components of the complexified subalgebra take the forms
\begin{equation}\label{eq_general_plane_now_element}
    S(\tau), \quad V(\tau)\mathbf{w}_{\perp}, \quad M(\tau)\mathbf{i}, \quad S,V,M \in \mathbb{C}.
\end{equation}
Because scalar contraction vanishes identically and the vector- and bivector-grade components lie entirely within the Plane-of-Now,
\begin{equation}
    \mathbf{U}\cdot S = 0, \quad \mathbf{U}\cdot(V\mathbf{w}_{\perp}) = 0, \quad \mathbf{U}\cdot(M\mathbf{i}) = 0.
\end{equation}
Each component therefore solves the algebraic propagation-orthogonality problem in Eq.~\eqref{eq_R_field_constraint} independently. Applying the temporal-volume admissibility condition $\mathbf{I}_{t}\mathbf{R}^{\dagger}\mathbf{I}_{t}=-\mathbf{R}$ to each component gives
\begin{equation}
    S^{*} = S, \quad V^{*} = V, \quad M^{*} = -M,
\end{equation}
with $^{*}$ denoting complex conjugation. The scalar- and vector-grade coefficients are therefore real, whereas the bivector-grade coefficient must be purely imaginary. Writing
\begin{equation}
    S = \mu(\tau), \quad V = g_{\pi}\varrho(\tau), \quad M = -iy\phi(\tau),
\end{equation}
where $y$ and $g_{\pi}$ are real constants and $\phi(\tau)$, $\varrho(\tau)$, and $\mu(\tau)$ are real-valued amplitudes evaluated along the trajectory, gives
\begin{equation}\label{eq_admissible_plane_now_source}
    \mathbf{R}_{\mathrm{PoN}} = \mu - iy\phi\mathbf{i} + g_{\pi}\varrho\mathbf{w}_{\perp}.
\end{equation}

\subsubsection{Even Grade Sector--Preserving Solutions: Mass and Scalar Yukawa Interactions}
Using the direct representation $R=\rho(\mathbf{R})$, the scalar-grade component yields
\begin{equation}
    R_{\mu} = \rho(\mu) = \mu I_{8},
\end{equation}
so a constant scalar component $\mu = mc$ supplies the mass term, whereas the bivector-grade component gives
\begin{equation}
    R_{\phi} = \rho(-iy\phi\mathbf{i}) = y\phi\Sigma_{z}.
\end{equation}
Their combined contribution to the PO-source interaction is
\begin{equation}\label{eq_scalar_yukawa_interaction}
    -\kappa\bar{\Psi}\Psi - \chi y\phi\bar{\Psi}\Sigma_{z}\Psi, \quad \kappa \coloneqq \frac{mc}{\hbar},
\end{equation}
where
\begin{equation}
    \bar{\Psi}\Psi = \bar{\Psi}_{+}\Psi_{+} - \bar{\Psi}_{-}\Psi_{-}, \quad \bar{\Psi}\Sigma_{z}\Psi = \bar{\Psi}_{+}\Psi_{+} + \bar{\Psi}_{-}\Psi_{-}.
\end{equation}
Thus, the scalar-grade component generates the mass term, while the Plane-of-Now bivector with a purely imaginary coefficient produces the scalar Yukawa interaction.

Adding the scalar-field contribution
\begin{equation}
    \mathcal{L}_{\phi} = \frac{1}{2}\partial_{\mu}\phi\partial^{\mu}\phi - \mathcal{V}_{\phi}(\phi)
\end{equation}
provides a dynamical realization of its propagation and self-interaction.

\subsubsection{Odd-Grade Sector--Mixing Solutions: Charged Pseudoscalar Fields and Isospin Exchange}
For the odd-grade sector, the vector component of the Plane-of-Now representation is
\begin{equation}\label{eq_plane_now_vector_operator}
    R_{\pi} = \rho(g_{\pi}\varrho\mathbf{w}_{\perp}) = g_{\pi}\varrho\rho(\mathbf{w}_{\perp}), \quad \mathbf{w}_{\perp} \coloneqq \mathbf{i}\mathbf{w},
\end{equation}
with
\begin{equation}
    \mathbf{w}_{\perp} = \sin\vartheta\mathbf{e}_{5} - \cos\vartheta\mathbf{e}_{6}.
\end{equation}
Using $\rho(\mathbf{e}_{5}) = \sigma_{x} \otimes \gamma_{5}$ and $\rho(\mathbf{e}_{6}) = \sigma_{y} \otimes \gamma_{5}$, the represented operator becomes
\begin{equation}\label{eq_plane_now_vector_explicit}
    R_{\pi} = g_{\pi}\varrho\left(\sin\vartheta\sigma_{x} - \cos\vartheta\sigma_{y}\right) \otimes \gamma_{5}.
\end{equation}
Defining the conjugate mixing fields by
\begin{equation}\label{eq_complex_pion_fields}
    \pi^{+}(\tau) \coloneqq \varrho e^{-i\vartheta}, \quad \pi^{-}(\tau) \coloneqq \varrho e^{i\vartheta} = \left(\pi^{+}\right)^{*},
\end{equation}
the represented PO-source operator takes the form
\begin{equation}\label{eq_plane_now_vector_complex}
    R_{\pi} = g_{\pi}\begin{pmatrix}0 & i\pi^{+}\gamma_{5} \\ -i\pi^{-}\gamma_{5} & 0\end{pmatrix}.
\end{equation}
Its contribution to the interaction is
\begin{equation}\label{eq_pion_interaction}
    -\chi\bar{\Psi}R_{\pi}\Psi = -i\chi g_{\pi}\left(\pi^{+}\bar{\Psi}_{+}\gamma_{5}\Psi_{-} + \pi^{-}\bar{\Psi}_{-}\gamma_{5}\Psi_{+}\right).
\end{equation}
Although $g_{\pi}\varrho\mathbf{w}_{\perp}$ has vector grade within the Plane-of-Now, its represented action is a scalar under proper four-dimensional Lorentz transformations, while its off-diagonal form exchanges the two Flag-time orientation sectors. For equal intrinsic parity phases of the displayed spinor columns, $\Psi_{\pm}'(t,\mathbf{x})=\gamma^{0}\Psi_{\pm}(t,-\mathbf{x})$, the identity $\gamma^{0}\gamma_{5}\gamma^{0}=-\gamma_{5}$ makes each mixing bilinear change sign, so the parity-covariant interaction has
\begin{equation}
    \pi^{\pm\prime}(t,\mathbf{x})=-\pi^{\pm}(t,-\mathbf{x}),
\end{equation}
giving the mixing fields their pseudoscalar character in this realization.

Including the mass and scalar Yukawa components, the field equations become
\begin{subequations}\label{eq_pion_coupled_sectors}
    \begin{align}
        \left(i\gamma^{\mu}D_{\mu}^{+} - \kappa - \chi y\phi\right)\Psi_{+} - i\chi g_{\pi}\pi^{+}\gamma_{5}\Psi_{-} &= 0,\\
        \left(i\gamma^{\mu}D_{\mu}^{-} + \kappa - \chi y\phi\right)\Psi_{-} - i\chi g_{\pi}\pi^{-}\gamma_{5}\Psi_{+} &= 0.
    \end{align}
\end{subequations}
For $\pi^{\pm}\neq0$, the two spinor sectors form a coupled eight-component field, whereas for $\pi^{\pm}=0$ they decouple into Dirac equations governed by opposite gauge phases. In this decoupled case with $\phi=0$, the transformation $\Psi_{-}\mapsto\gamma_{5}\Psi_{-}$ brings the lower equation to the same positive-mass convention as the upper equation. A charge-conjugate pair is then given by $\Psi_{-}=\gamma_{5}C\bar{\Psi}_{+}^{\mathsf T}$, where $C$ is the charge-conjugation matrix satisfying $C\gamma^{\mu\mathsf T}C^{-1}=-\gamma^{\mu}$.

For this pairing, the charge-conjugate Dirac field has opposite intrinsic parity, while the additional $\gamma_{5}$ reverses that sign, yielding the equal displayed parity phases used above. With a nonzero scalar background, the effective masses become $mc+y\phi$ and $mc-y\phi$ after the lower-block transformation, so charge conjugation connects the corresponding backgrounds through $\phi\mapsto-\phi$.

Reversal of the Pole-time propagation branch exchanges the eigenvalues $+c \leftrightarrow -c$, whereas the charge-conjugate realization relates the two Flag-time charge orientations. Their opposite charges therefore remain distinct from the choice of future-directed causal evolution.

Introduce the transition and orientation operators
\begin{equation}
    \mathcal{I}_{\pm} \coloneqq \frac{1}{2}\left(\sigma_{x} \pm i\sigma_{y}\right) \otimes \gamma_{5}, \quad \mathcal{I}_{z} \coloneqq \frac{1}{2}\Sigma_{z}.
\end{equation}
They satisfy
\begin{equation}\label{eq_isospin_ladder_algebra}
    [\mathcal{I}_{z}, \mathcal{I}_{\pm}] = \pm\mathcal{I}_{\pm}, \quad [\mathcal{I}_{+}, \mathcal{I}_{-}] = 2\mathcal{I}_{z},
\end{equation}
and form a Plane-of-Now $\mathfrak{su}(2)$ ladder algebra. Since the vector-grade source also obeys $\{U,R_{\pi}\}=0$, its action exchanges propagation eigenspaces, with this exchange balanced by the derivative terms in propagation-aligned solutions. In this basis, the mixing operator becomes
\begin{equation}\label{eq_pion_ladder_decomposition}
    R_{\pi} = ig_{\pi}\pi^{+}\mathcal{I}_{+} - ig_{\pi}\pi^{-}\mathcal{I}_{-},
\end{equation}
where
\begin{equation}
    \mathcal{I}_{+} : \Psi_{-} \to \Psi_{+}, \quad \mathcal{I}_{-} : \Psi_{+} \to \Psi_{-}.
\end{equation}
Under the transformation
\begin{equation}
    \Omega(\alpha) = \exp\left(-i\alpha\mathcal{I}_{z}\right),
\end{equation}
the transition operators transform as
\begin{equation}
    \Omega\mathcal{I}_{\pm}\Omega^{-1} = e^{\mp i\alpha}\mathcal{I}_{\pm}.
\end{equation}
Covariance of $R_{\pi}$ therefore gives
\begin{equation}
    \pi^{\pm} \to e^{\mp i\alpha}\pi^{\pm},
\end{equation}
so the conjugate mixing fields have weights $\pm1$ under the $\mathrm{U}(1)$ subgroup generated by $\mathcal{I}_{z}$, while the spinor columns transform as $\Psi_{\pm}\to e^{\mp i\alpha/2}\Psi_{\pm}$ and have weights $\pm\frac{1}{2}$. Under local changes of this orientation reference, the Principle of Relativity carries these weights into the covariant derivatives,
\begin{equation}\label{eq_pion_gauge_covariant_derivatives}
    D_{\mu}\pi^{\pm} = \left(\partial_{\mu} \pm \frac{i}{\hbar}q_{\mathrm{I}}B_{\mu}^{\mathrm{I}}\right)\pi^{\pm},
\end{equation}
with the corresponding spinor derivatives
\begin{equation}
    D_{\mu}\Psi_{\pm}=\left(\partial_{\mu}\pm\frac{iq_{\mathrm I}}{2\hbar}B_{\mu}^{\mathrm I}\right)\Psi_{\pm}.
\end{equation}
Thus, the mixing-field charge is the difference between the two spinor charges, giving $\pm2q$ when the spinor charges are $\pm q$. If this subgroup is identified with the electromagnetic coupling, $q_{\mathrm I}B_{\mu}^{\mathrm I}=2qA_{\mu}$, consistently reproducing both derivatives. For a nonzero mixing operator,
\begin{equation}
    [\mathcal{I}_{z}, R_{\pi}] = ig_{\pi}\pi^{+}\mathcal{I}_{+} + ig_{\pi}\pi^{-}\mathcal{I}_{-} \neq 0,
\end{equation}
so the isospin current
\begin{equation}
    J_{\mathrm{I}}^{\mu} \coloneqq \bar{\Psi}\Gamma^{\mu}\mathcal{I}_{z}\Psi
\end{equation}
satisfies
\begin{equation}\label{eq_pion_isospin_exchange}
    \partial_{\mu}J_{\mathrm{I}}^{\mu} = -i\chi\bar{\Psi}[\mathcal{I}_{z}, R_{\pi}]\Psi.
\end{equation}
Consequently, the mixing field carries the charge exchanged between the two spinor sectors.

Completing the charged-field dynamics with
\begin{equation}
    \mathcal{L}_{\pi} = D_{\mu}\pi^{-}D^{\mu}\pi^{+} - \mathcal{V}_{\pi}\left(\pi^{-}\pi^{+}\right)
\end{equation}
yields the corresponding current
\begin{equation}
    j_{\pi}^{\mu} = i\left(\pi^{-}D^{\mu}\pi^{+} - D^{\mu}\pi^{-}\pi^{+}\right).
\end{equation}
Its field equations give the complementary exchange balance
\begin{equation}
    \partial_{\mu}j_{\pi}^{\mu} = i\chi\bar{\Psi}[\mathcal{I}_{z}, R_{\pi}]\Psi,
\end{equation}
so adding the two balances gives
\begin{equation}\label{eq_total_pion_isospin_conservation}
    \partial_{\mu}\left(J_{\mathrm{I}}^{\mu} + j_{\pi}^{\mu}\right) = 0.
\end{equation}
Together, these relations give the Plane-of-Now vector solution a pion-type charged-pseudoscalar interaction, with isospin raising and lowering operators that exchange the two spinor sectors and a conserved total orientation charge.

\subsection{The Interacting Field Lagrangian}
Including the scalar mass contribution within $\mathbf{R}$, the complete PO-source is
\begin{equation}
    \mathbf{R}_{\mathrm{total}} \equiv \mathbf{R}= \mathbf{R}_{\mathrm{PoN}} + \mathbf{R}_{\mathrm{B}} + \mathbf{R}_{\mathrm{Other}},
\end{equation}
where $\mathbf{R}_{\mathrm{B}}=\mathbf{i}qa\left(|B|^{2}\right)\mathbf{B}$ is the propagation-orthogonal field-strength solution derived previously, with $\rho(\mathbf{R}_{\mathrm{B}})=R_{\mathrm{B}}=iqa\Sigma_{z}B$, while $\mathbf{R}_{\mathrm{Other}}$ contains the remaining admissible components of the general six-dimensional algebra established in Appendix~\ref{appx_general_matrix_source}.

Combining these contributions gives the represented PO-source operator
\begin{equation}\label{eq_full_geometric_source_operator}
\begin{split}
    R ={}& mcI_{8} + y\phi\Sigma_{z} + ig_{\pi}\left(\pi^{+}\mathcal{I}_{+} - \pi^{-}\mathcal{I}_{-}\right) \\
    & +iqa\Sigma_{z}B+\rho\left(\mathbf{R}_{\mathrm{Other}}\right).
\end{split}
\end{equation}
As established in Section~\ref{ch_gauge_covariance}, local reference covariance follows from the Principle of Relativity, so the derivative enters as $D_{\mu}\coloneqq\partial_{\mu}+\frac{i}{\hbar}\sum_{a}q_{a}A_{\mu}^{a}T_{a}$, using Hermitian interaction generators $T_{a}$ in this convention. Consequently, the kinetic-derivative term provides the corresponding minimal gauge interactions.

Collecting the kinetic, source, and field contributions gives the full interacting Lagrangian,
\begin{equation}\label{eq_full_field_lagrangian}
\begin{split}
    \mathcal{L} ={}& \frac{i}{2}\bar{\Psi}\Gamma^{\mu}\overleftrightarrow{D}_{\mu}\Psi - \chi\bar{\Psi}R\Psi \\
    & + \mathcal{L}_{\mathrm{F}} + \mathcal{L}_{\mathrm{B}} + \mathcal{L}_{\phi} + \mathcal{L}_{\pi} + \mathcal{L}_{\mathrm{R}}^{\mathrm{Other}}.
\end{split}
\end{equation}
Here the covariant derivative carries the gauge interactions of the chosen field representation, with generators commuting with the kinetic matrices and sources transforming as $R^{\prime}=gRg^{-1}$. Through its admissible components, the PO-source supplies the spinor mass, scalar Yukawa, charged-pseudoscalar mixing, propagation-orthogonal field-strength, and remaining algebraic interactions, while the accompanying dynamical terms govern the source fields. Specifying the representations, fields, potentials, and coupling constants then selects a particular physical realization.

Finally, the propagation-orthogonal field-strength coupling $-\chi\bar{\Psi}R_{\mathrm{B}}\Psi$ is conventionally classified as a nonminimal electromagnetic field-strength interaction. In the present theory, however, such terms belong to the minimal interaction structure, arising directly from the geometric-source equation alongside the usual minimal electromagnetic coupling. Accordingly, anomalous magnetic responses and nonlinear electromagnetic couplings emerge as expected solutions of the theory and not as ad hoc additions.
